%% End matter, in the order Elsevier's guide gives: CRediT, competing interests,
%% data availability, acknowledgements (with funding), generative-AI declaration.
%% TODO(author): confirm the CRediT roles and the funding sentence before submitting.

\section*{CRediT authorship contribution statement}

\textbf{Nambi Rajan M:} Conceptualization, Methodology, Software, Validation, Formal
analysis, Investigation, Data curation, Writing -- original draft, Writing -- review \&
editing, Visualization.

\section*{Declaration of competing interest}

The author declares that they have no known competing financial interests or personal
relationships that could have appeared to influence the work reported in this paper.

\section*{Data availability}

The compiler, the Coq development, the benchmark suites, the real-workflow corpus with
its protocol, exclusion log and labels, and the evaluation harness are openly available at
\url{https://github.com/guyoverclocked/orchlang-compiler-design-project}. The commands
\code{make check}, \code{make proofs} and \code{python bench/evaluate.py} reproduce every
number reported in this paper, and \file{submission/scp/figures/data.py} regenerates the
data behind every plot.
%% TODO(author): archive a tagged release on Zenodo and add its DOI here; add a licence.

\section*{Acknowledgements}

This research did not receive any specific grant from funding agencies in the public,
commercial, or not-for-profit sectors.

\section*{Declaration of generative AI and AI-assisted technologies in the manuscript preparation process}

During the preparation of this work the author used Claude (Anthropic) in order to
develop parts of the compiler, the evaluation harness and the Coq proofs (as described in
\cref{sec:implementation}), to draw the explanatory diagrams, and to draft and edit the
text and check the references against publisher records. After using this tool, the
author reviewed and edited the content as needed and takes full responsibility for the
content of the published article.
